% ------------------------------------------------------------------
%  Appendice A — Una proposta di indagine: come studiano gli studenti italiani
%  Ricerca di riferimento: ricerca 05 (progetto di indagine, v1.0)
%  Scheda nell'indice ragionato: ricerca/06_indice-ragionato.md, §6.3
%    (decisione 6: nessun risultato; l'indagine si farà con un partner,
%    dopo la pubblicazione)
%  Voci bibliografiche: karpicke2009, kornell2007, hartwig2012,
%    mccabe2011, roediger2006, mcdaniel2020, dekker2012, bei2023
%
%  STATO: prima stesura (07/10/2026), circa 3.300 parole compreso il
%  nucleo del questionario (pagine finali, fotocopiabili).
%
%  PROMESSE DA MANTENERE (rimandi in entrata):
%  - cap. 2, esercizio: «elenca le strategie che usi … in ordine di
%    frequenza … accanto a ciascuna scrivi perché la usi. È la stessa
%    domanda che l'Appendice propone»: domanda 1 del nucleo (A.8),
%    formulata con le stesse parole;
%  - cap. 2, testo: «l'Appendice propone un modo per farlo» (A.1–A.7);
%  - cap. 2, riquadro «Per chi insegna»: «per un questionario più
%    completo si può partire da quello proposto nell'Appendice» (A.7, A.8);
%  - cap. 8, «E in Italia?»: «Nell'appendice A propongo come farlo
%    [l'esperimento di Roediger e Karpicke in un'aula italiana], e invito
%    le università e le scuole» (A.5, A.7).
%
%  DATI: nessun dato nuovo sugli studenti italiani. Dati già verificati:
%  Kornell & Bjork 2007 e Hartwig & Dunlosky 2012 (registro S-29),
%  Karpicke et al. 2009 (S-31, solo richiamato), Roediger & Karpicke 2006
%  (cap. 8, solo richiamato). McCabe 2011: verificato il 07/10/2026
%  sull'abstract e sulla pagina dell'editore (255 studenti, accuratezza
%  media 23%, 71% nel seminario che aveva letto gli studi originali;
%  registro S-29, ora chiuso anche per McCabe). Calcoli propri (in nota):
%  campione per ±3 e ±5 punti (1.067 e 385); potenza dell'esperimento
%  entro i soggetti (d = 0,4, alfa 0,05 bilaterale, potenza 0,80:
%  52 studenti; d = 0,5: 34). Stime di tempo dalla ricerca 05 §3.3,
%  presentate come stime.
%  Voce: «io» discreto (propongo); l'autore non si presenta come
%  ricercatore né come docente (decisione 13). Nel nucleo del
%  questionario «tu» è lo studente che risponde.
%  Nomi negli scenari: Giulia e Marco (come in ricerca 05); «Luca» della
%  ricerca 05 sostituito con «Matteo».
% ------------------------------------------------------------------
\chapter{Una proposta di indagine: come studiano gli studenti italiani}
\label{app:indagine}

%  Comandi usati solo in questa appendice (nomi diversi da quelli
%  dell'Appendice B, che li definisce per conto suo)
\newcommand{\quadretto}{\tikz[baseline=-0.05em]\draw[line width=0.4pt]
  (0,0) rectangle (0.62em,0.62em);}
\newcommand{\rigavuota}{\rule{0pt}{1.6\baselineskip}}
%  \domanda{n.}: nuova domanda; se in fondo alla pagina restano meno di
%  sei righe, la domanda passa alla pagina seguente (niente domande
%  separate dalle loro risposte)
\newcommand{\domanda}[1]{\par\medskip
  \begingroup\dimen0=\pagegoal\advance\dimen0 by -\pagetotal
  \ifdim\dimen0<6\baselineskip\newpage\fi\endgroup
  \noindent\textbf{#1}\enspace}
\newenvironment{opzioni}
  {\begin{itemize}[label=\quadretto, leftmargin=1.6em, itemsep=0pt, topsep=0.2\baselineskip]}
  {\end{itemize}}

\iniziale{N}{el capitolo}~\ref{cap:oggi} ti ho chiesto di elencare le
strategie che usi quando studi per un esame, in ordine di frequenza, e
di scrivere accanto a ciascuna perché la usi. Se hai fatto l'esercizio,
hai risposto alla prima domanda di un'indagine che in Italia nessuno ha
ancora condotto. Negli Stati Uniti la stessa domanda è stata posta a
centinaia di studenti universitari, e le risposte le hai lette nel
capitolo~\ref{cap:vita}: quasi tutti rileggono, quasi nessuno si
interroga da solo per imparare \parencite{karpicke2009}. Degli
studenti italiani, invece, non sappiamo quanti rileggono, quanti si
mettono alla prova, quanti hanno avuto qualcuno che spiegasse loro come
si studia.

Questa appendice non contiene risultati. Contiene una proposta: che
cosa chiedere, a chi, come, con quali garanzie. È scritta per chi
potrebbe realizzarla -- gruppi di ricerca di psicologia dell'educazione
e di pedagogia sperimentale, centri per la didattica degli atenei,
servizi di orientamento e di tutorato, associazioni studentesche -- e
per chi vuole soltanto usarne le domande nella propria aula. Un'indagine
di questo tipo non si improvvisa e non si fa da soli: servono un gruppo
di ricerca, un comitato etico, atenei e scuole disposti a partecipare.
Per questo il libro la propone e la lascia a chi vorrà raccoglierla,
con un invito che trovi alla fine.

\section{Perché serve}

Le indagini su come studiano gli universitari, raccontate nei
capitoli~\ref{cap:vita} e~\ref{cap:oggi}, vengono quasi tutte dagli
Stati Uniti. Karpicke e colleghi hanno chiesto agli studenti di elencare
le proprie strategie \parencite{karpicke2009}; Kornell e Bjork hanno
chiesto a 472 studenti come decidevano che cosa studiare e perché si
interrogavano \parencite{kornell2007}; Hartwig e Dunlosky hanno messo
in relazione le abitudini di 324 studenti con la loro media dei voti
\parencite{hartwig2012}. Jennifer McCabe ha fatto un passo diverso: ha
descritto a 255 studenti sei situazioni di studio, ciascuna basata su un
esperimento noto, e ha chiesto quale dei due studenti descritti avrebbe
imparato di più. In media gli studenti hanno indovinato il 23\% delle
situazioni; quelli che avevano letto gli esperimenti originali in un
seminario, il 71\% \parencite{mccabe2011}. Non sapere come si impara non
è dunque un destino: è una lacuna, e si può colmare.

Non c'è motivo di pensare che la memoria di uno studente italiano
funzioni in modo diverso da quella di uno studente americano. Ma il
modo di studiare dipende anche dalle regole del gioco, e quelle in
Italia sono diverse: l'esame orale, la sessione che concentra più esami
in poche settimane, gli appelli che si possono ripetere, il voto che si
può rifiutare, le trascrizioni delle lezioni che passano di anno in
anno, il manuale come unica fonte, le ripetizioni private alle superiori.
E da qualche anno c'è un'abitudine nuova, l'intelligenza artificiale,
che nel capitolo~\ref{cap:oggi} abbiamo visto già diffusa tra i
quindicenni. Che cosa produca tutto questo nel modo di studiare degli
italiani, oggi, si può soltanto immaginare.

In Italia esistono questionari validati per valutare il metodo di
studio di un singolo studente, usati da psicologi e insegnanti.%
\footnote{Per esempio la batteria AMOS e gli altri strumenti di Cesare
  Cornoldi, Rossana De Beni e del Gruppo MT, pubblicati da Erickson
  \parencite{cornoldi}. Sono coperti da diritti d'autore: non si possono
  riprodurre in un'altra indagine senza un accordo con l'editore.}
Servono
però a un'altra cosa: a descrivere una persona per aiutarla, non a
fotografare un Paese e a confrontarlo con le ricerche internazionali.
Ciò che manca è un quadro nazionale, costruito con le stesse domande
delle indagini straniere, così da poter dire non soltanto \enquote{che
cosa fanno gli studenti italiani}, ma anche \enquote{se fanno come gli
altri, o diversamente}.

Un principio, prima di tutto il resto: meglio un'indagine più piccola
ma pulita -- domande provate, campione descritto con onestà, analisi
decise in anticipo -- che un grande sondaggio diffuso sui social, a cui
rispondono soltanto gli studenti già interessati al tema.

\section{Le domande}

L'indagine dovrebbe rispondere a nove domande, le prime tre riprese
direttamente dalle indagini straniere, per poter confrontare i
risultati.

\begin{enumerate}[leftmargin=1.6em, itemsep=0.2\baselineskip]
  \item Quali strategie usano più spesso gli studenti italiani, e quale
        indicano come principale? Quanti rileggono, quanti si
        interrogano da soli?
  \item Chi si interroga da solo, lo fa per \emph{imparare} o per
        \emph{controllare} quanto sa?
  \item Come distribuiscono lo studio nel tempo? Quanta parte si
        concentra nelle ultime settimane, o negli ultimi giorni, prima
        dell'appello?
  \item Che cosa \emph{credono} su come si impara? Se si descrivono
        loro due modi di studiare -- rileggere o provare a ricordare,
        studiare tutto in un giorno o in più giorni, esercizi a blocchi
        o mescolati -- quale scelgono, e con quanta sicurezza?
  \item Qualcuno ha mai insegnato loro come si studia? Chi: un
        insegnante, un corso, la famiglia, internet, un'intelligenza
        artificiale?
  \item Come usano l'intelligenza artificiale: per farsi spiegare e
        interrogare, o per farsi fare riassunti, testi ed esercizi?
  \item Come gestiscono il telefono mentre studiano? Quanto dormono
        prima di un esame?
  \item Come cambia tutto questo tra aree disciplinari, anni di corso,
        tipi di ateneo e regioni; tra chi è il primo della famiglia a
        frequentare l'università e chi no; tra gli studenti con un
        disturbo specifico dell'apprendimento e gli altri?
  \item Le strategie dichiarate sono associate ai risultati -- la media
        dei voti, gli esami in regola con il piano di studi -- a parità
        delle altre condizioni?
\end{enumerate}

L'ultima domanda va presa con cautela. Un'indagine con questionario
mostra associazioni, non cause: se chi si interroga da solo ha voti più
alti, può darsi che lo debba al metodo, ma può darsi anche che studi di
più, o che parta avvantaggiato. Ciò che una strategia \emph{produce}
lo dicono gli esperimenti, come quelli raccontati nella terza parte del
libro, non i sondaggi.

Un breve modulo per chi insegna potrebbe affiancare quello per gli
studenti: quali strategie consigliano i docenti, quali ritengono
efficaci, quanti conoscono la pratica del recupero e la distribuzione,
quanti credono ai miti di cui parla il capitolo~\ref{cap:miti}, primo
fra tutti quello degli stili di apprendimento
\parencite{dekker2012,bei2023}. Confrontare le risposte dei docenti con
quelle degli studenti direbbe molto su dove nasce il metodo che gli
studenti si portano dietro.

Le indagini straniere permettono di prevedere alcune risposte: che la
rilettura sarà la strategia più citata, che chi si interroga lo farà
soprattutto per controllarsi, che pochi diranno di aver ricevuto un
insegnamento esplicito sul metodo. Sono ipotesi, non risultati, e vanno
scritte e depositate prima di raccogliere i dati (lo vedremo nel
paragrafo~\ref{sec:indagine-garanzie}). Se i dati italiani le
smentissero, sarebbe la notizia più interessante dell'indagine.

\section{Il questionario}
\label{sec:indagine-questionario}

Il cuore dell'indagine è un questionario anonimo, da compilare online in
quindici o venti minuti. Le pagine finali di questa appendice ne
riportano il nucleo, nella versione per l'università; qui ne descrivo
la struttura e le regole con cui è costruito.

\begin{table}[tbp]\small
  \centering
  \begin{tabular}{@{}>{\raggedright\arraybackslash}p{0.22\textwidth}>{\raggedright\arraybackslash}p{0.41\textwidth}>{\raggedright\arraybackslash}p{0.27\textwidth}@{}}
    \toprule
    \textit{Parte} & \textit{Che cosa chiede} & \textit{Modello} \\
    \midrule
    Le tue strategie & elenco libero, in ordine di frequenza, con il perché; la strategia più utile & \textcite{karpicke2009} \\
    Una scelta & dopo aver letto un capitolo: rileggere o provare a ricordare? E perché? & \textcite{karpicke2009} \\
    Che cosa fai & frequenza di una quindicina di strategie; perché ci si interroga; che cosa si fa dopo l'esame & \textcite{kornell2007}; \textcite{hartwig2012} \\
    Chi ti ha insegnato & se e da chi si è imparato a studiare & \textcite{kornell2007}, ampliato \\
    Il tempo & distribuzione dello studio prima dell'ultimo esame; esami per sessione; sonno & \textcite{hartwig2012} \\
    Che cosa pensi & situazioni di studio a scelta, con il grado di sicurezza & \textcite{mccabe2011} \\
    Intelligenza artificiale & quanto e per che cosa & originale \\
    Attenzione & telefono, notifiche, durata delle sessioni & originale \\
    Esami e motivazione & prove intermedie, ricordo a sei mesi, perché si studia & originale \\
    Su di te & area e anno di corso, ateneo, regione, titolo di studio dei genitori, media dei voti; facoltativa, la certificazione di DSA & originale \\
    \bottomrule
  \end{tabular}
\end{table}

Contano, più delle singole domande, alcune regole.

\emph{Prima le domande aperte, poi quelle chiuse.} Se lo studente vede
subito un elenco di strategie, ne riconosce alcune e le dichiara anche
se non le usa. Per questo l'elenco libero viene per primo, prima di
qualunque suggerimento.

\emph{Le opinioni dopo i comportamenti.} Le situazioni a scelta sono
istruttive: chi le legge comincia a sospettare quale sia la risposta
giusta, e potrebbe correggere di conseguenza ciò che dichiara di fare.
Vanno quindi dopo tutte le domande su ciò che lo studente fa davvero.

\emph{Nessuna risposta giusta in vista.} Nessuna domanda deve far
capire quale strategia sia preferibile; l'ordine delle strategie e delle
situazioni viene mescolato a caso da uno studente all'altro, e la prima
schermata lo dice chiaramente: interessa ciò che fai, non ciò che
\enquote{si dovrebbe} fare.

\emph{Domande concrete.} \enquote{L'ultimo esame che hai sostenuto} è
meglio di \enquote{di solito}: la memoria di un episodio preciso è meno
esposta al desiderio di fare bella figura.

\emph{Controlli di qualità.} Una riga di attenzione (\enquote{per
questa riga scegli \enquote{Spesso}}) e un tempo minimo di compilazione
permettono di scartare chi risponde a caso; i criteri si fissano prima
di vedere i dati.

\emph{Una versione per ogni livello.} Il nucleo resta lo stesso, per
permettere i confronti; cambiano le parole. Per la scuola si parla di
verifiche e interrogazioni, e alle medie le frasi si fanno più brevi e
gli esempi più concreti.

Le risposte all'elenco libero vanno poi classificate -- rilettura,
sottolineatura, riassunto, ripetizione a memoria, esercizi, mappe e
così via -- da due persone che lavorano separatamente, con un calcolo
del loro accordo e una discussione dei casi dubbi. Un'intelligenza
artificiale può fare da terzo controllo, mai da sola: per le ragioni
del capitolo~\ref{cap:ia}, il giudizio su che cosa uno studente intende
dire spetta a chi può risponderne.

\section{A chi chiedere}

Il libro, in questa prima edizione, parla soprattutto agli studenti
universitari, e da loro conviene cominciare. In un secondo momento lo
stesso questionario, adattato, potrebbe essere proposto agli studenti
delle superiori e delle medie (la primaria è esclusa: i bambini non
sanno ancora descrivere in modo affidabile come studiano).

Quanti studenti servono? Dipende da che cosa si vuole sapere. Per
stimare una percentuale nazionale con un margine di circa tre punti
servono poco più di mille risposte; per confrontare gruppi diversi --
le aree disciplinari, per esempio -- ne servono circa quattrocento per
ciascun gruppo.\footnote{Sono i valori classici per un campione
casuale, con un livello di confidenza del 95\% e nel caso più
sfavorevole (una percentuale vicina al 50\%): circa 1.070 risposte per
un margine di ±3 punti, circa 385 per ±5. Con cinque aree
disciplinari si arriva a circa duemila studenti.} Un obiettivo
realistico è di qualche migliaio di studenti e di qualche centinaio di
docenti.

C'è però una cautela che nessun numero risolve. Un'indagine di questo
tipo raccoglie quasi sempre un campione \emph{di convenienza}: rispondono
gli atenei e i corsi che accettano di partecipare. Per ridurre le
distorsioni conviene somministrare il questionario in aula, durante una
lezione concordata con il docente, invece di affidarsi a un link sui
social; coinvolgere atenei grandi e piccoli, statali e non statali, di
tutte le regioni, e corsi di tutte le aree; confrontare la composizione
del campione con i dati nazionali sugli iscritti e, se serve,
ricalibrare le stime. E, soprattutto, dichiarare i limiti: un campione
di convenienza descritto con onestà vale più di un campione che si
spaccia per rappresentativo.

\section{Un esperimento in aula}
\label{sec:indagine-esperimento}

Un questionario dice che cosa gli studenti \emph{credono}; un
esperimento mostra che cosa \emph{funziona}. Nel
capitolo~\ref{cap:recupero} ho raccontato l'esperimento di Roediger e
Karpicke: chi rileggeva era più sicuro di sé, ma una settimana dopo
ricordava meno di chi aveva chiuso il libro e provato a ricordare
\parencite{roediger2006}. Nessuno, finora, lo ha ripetuto in un'aula
italiana. Propongo di farlo, accanto all'indagine, in una forma
semplice.

\begin{enumerate}[leftmargin=1.6em, itemsep=0.2\baselineskip]
  \item \emph{Il materiale.} Due brevi testi divulgativi, di
        duecentocinquanta o trecento parole ciascuno, su argomenti che
        non fanno parte del corso (per evitare che qualcuno li conosca
        già), di lunghezza e difficoltà simili, controllate con una
        prova preliminare.
  \item \emph{Le due condizioni.} Ogni studente legge entrambi i testi
        per cinque minuti. Poi, su uno dei due, passa altri cinque
        minuti a rileggerlo; sull'altro, cinque minuti a scrivere a
        libro chiuso tutto ciò che ricorda. Quale testo vada con quale
        condizione, e in che ordine, si decide a caso: così ognuno è
        il termine di confronto di se stesso.
  \item \emph{La previsione.} Subito dopo, ogni studente indica quanto
        pensa che ricorderà di ciascun testo tra una settimana, da 0 a
        100.
  \item \emph{La prova.} Una settimana dopo, senza preavviso, domande a
        risposta breve su entrambi i testi, corrette da due persone in
        modo indipendente.
  \item \emph{La restituzione.} Nella stessa lezione, o in quella
        successiva, gli studenti vedono i risultati del proprio gruppo
        -- quanto ricordavano dei testi riletti e di quelli richiamati,
        e quanto pensavano di ricordarne -- e qualcuno spiega loro che
        cosa significano.
\end{enumerate}

Non servono numeri enormi: se l'effetto è di grandezza media, come
nelle ricerche straniere, una cinquantina di studenti basta per
vederlo.\footnote{Calcolo per un confronto entro gli stessi soggetti,
con una probabilità di errore del 5\% e una potenza dell'80\%: 52
studenti per un effetto di grandezza \emph{d}~=~0,4, 34 per \emph{d}~=~0,5.
Il calcolo va rifatto sul disegno definitivo.} Ma qualche classe in più,
in corsi e in atenei diversi, permetterebbe di vedere se l'effetto
cambia da una disciplina all'altra. Il disegno ammette varianti: una
terza condizione in cui, dopo aver provato a ricordare, si ricontrolla
il testo; oppure una versione sulla distribuzione, con lo stesso tempo
di studio in una seduta sola o in due sedute a qualche giorno di
distanza.

L'esperimento ha un valore che va oltre i dati. Una cosa è leggere che
la pratica del recupero funziona; un'altra è vederlo sui propri
risultati, accanto alla propria previsione sbagliata. Le ricerche
sull'insegnamento delle strategie di studio indicano proprio questo
passaggio -- convincersi, oltre che sapere -- come uno dei più
importanti perché un'abitudine cambi davvero \parencite{mcdaniel2020}.
Per gli studenti che partecipano, la restituzione è già una lezione sul
metodo.

\section{Le garanzie}
\label{sec:indagine-garanzie}

Un'indagine che chiede agli studenti come studiano maneggia dati
personali, a volte di minorenni, a volte relativi alla salute. Le
garanzie non sono un dettaglio burocratico: sono la condizione perché
le risposte siano sincere e perché i risultati siano credibili. Ecco
quelle essenziali. I documenti definitivi -- informativa, moduli di
consenso, eventuale valutazione d'impatto -- andranno preparati con chi
si occupa di protezione dei dati nell'ateneo partner.

\emph{Partecipazione libera e anonima.} Si può rifiutare o smettere in
qualunque momento, senza alcuna conseguenza su voti ed esami. Nessun
nome, indirizzo email o indirizzo IP; si raccolgono solo i dati che
servono alle domande dell'indagine.

\emph{Un comitato etico.} Il progetto va sottoposto al comitato etico
per la ricerca dell'ateneo partner, che è anche il soggetto più adatto
a essere titolare del trattamento dei dati. Se un giorno il questionario
arrivasse nelle scuole, servirebbero l'autorizzazione del dirigente, il
consenso dei genitori per i minorenni e l'assenso degli studenti stessi,
che devono poter dire di no.

\emph{Prudenza con i dati sulla salute.} La domanda su una
certificazione di disturbo specifico dell'apprendimento riguarda dati
relativi alla salute, che il Regolamento europeo sulla protezione dei
dati tutela in modo particolare (art.~9). Deve essere facoltativa, con
un consenso specifico; e nessun risultato va pubblicato per gruppi così
piccoli -- meno di dieci studenti, per esempio -- da rendere riconoscibile
qualcuno. Se questo non si può garantire, è meglio porre la domanda
soltanto agli studenti maggiorenni, o non porla.

\emph{Le ipotesi prima dei dati.} Domande, ipotesi, criteri di
esclusione e analisi vanno scritti e depositati in un archivio pubblico
prima di raccogliere i dati, come si fa con la preregistrazione, per
esempio sull'\emph{Open Science Framework}. È la garanzia che i
risultati non siano stati scelti dopo averli visti.

\emph{Tutto pubblico, tranne le persone.} Il questionario, il codice
delle analisi e i dati resi anonimi vanno pubblicati, perché chiunque
possa controllare e ripetere. Il questionario di queste pagine ha la
stessa licenza del libro: si può copiare e adattare, citando la fonte.

\emph{Un conflitto d'interessi dichiarato.} Chi propone l'indagine ha
scritto un libro che ne attende i risultati, e potrebbe sperare che gli
diano ragione. Va detto apertamente; ed è un motivo in più per
depositare le ipotesi in anticipo e pubblicare tutto, anche ciò che le
smentisce. Se i dati italiani contraddicessero qualcosa di ciò che si
legge in queste pagine, una nuova edizione del libro dovrebbe dirlo.

\emph{Qualcosa in cambio.} Agli atenei e ai corsi che partecipano
spetta un rapporto con i risultati aggregati; agli studenti, alla fine
del questionario, una pagina con i principi di studio che la ricerca
sostiene. Chi partecipa dovrebbe uscirne sapendo qualcosa in più.

\section{Come procedere}

Le tappe sono quelle di qualunque indagine seria.

\begin{enumerate}[leftmargin=1.6em, itemsep=0.2\baselineskip]
  \item \emph{Un partner.} Un gruppo di ricerca universitario che
        prenda la responsabilità scientifica ed etica del progetto.
  \item \emph{La revisione del questionario.} Due o tre esperti -- di
        psicologia dell'apprendimento, di metodologia della ricerca
        educativa, di didattica universitaria -- lo leggono e lo
        correggono.
  \item \emph{Le interviste cognitive.} Cinque-otto studenti per gruppo
        compilano il questionario ad alta voce, per verificare che ogni
        domanda sia capita come previsto.
  \item \emph{La prova pilota.} Qualche decina di studenti, per
        misurare i tempi, trovare le domande ambigue e controllare la
        piattaforma (che dovrebbe conservare i dati nell'Unione europea
        e poter mescolare le domande).
  \item \emph{La preregistrazione}, con la versione definitiva del
        questionario.
  \item \emph{La raccolta}, lontano dalle sessioni d'esame, insieme
        all'esperimento in aula.
  \item \emph{La codifica e l'analisi}, secondo il piano depositato.
  \item \emph{La restituzione}: rapporto agli atenei, dati e codice
        pubblici, un articolo scientifico se il partner lo ritiene
        opportuno.
\end{enumerate}

Una versione essenziale, con un campione di convenienza ben descritto,
richiede indicativamente da sei a nove mesi; una versione con un
campione costruito per essere rappresentativo e destinata a una rivista
scientifica, da dodici a diciotto. I costi sono modesti: il lavoro di
chi coordina, una piattaforma per i questionari, la consulenza sulla
protezione dei dati e qualche ora di codifica, che studenti in tirocinio
potrebbero svolgere imparando, nel frattempo, il mestiere della ricerca.

\section{Un invito}

Se lavori in un gruppo di ricerca, in un centro per la didattica, in un
servizio di orientamento o di tutorato, o rappresenti gli studenti in
un ateneo, e questa proposta ti interessa, scrivimi:
\texttt{luca@lucalevi.com}.
Metterò a disposizione tutto il materiale preparato per il libro, a
cominciare dalla versione completa del questionario, e lascerò al
gruppo di ricerca le scelte che spettano a chi ne porta la
responsabilità.

Se insegni e vuoi soltanto sapere come studiano i tuoi studenti, non
devi aspettare nessuno. Puoi usare le domande delle pagine seguenti
nella prima lezione del corso, su carta e senza nomi, come una normale
attività didattica: è il suggerimento del capitolo~\ref{cap:oggi}, in
forma più completa. Se un giorno volessi raccoglierne i risultati per
pubblicarli, valgono le garanzie del
paragrafo~\ref{sec:indagine-garanzie}; e il primo passo è parlarne con
il comitato etico del tuo ateneo. Le risposte più utili alle domande
sulle situazioni di studio sono nei capitoli~\ref{cap:recupero},
\ref{cap:tempo} e~\ref{cap:alternanza}: è bene discuterle con gli
studenti \emph{dopo} che hanno risposto, non prima.

E se sei uno studente e hai conservato il foglio del
capitolo~\ref{cap:oggi}, riprendilo adesso e confrontalo con la
domanda~1 delle pagine che seguono. È la stessa. Le tue risposte, oggi,
sono ancora quelle di allora?

%% =================================================================
%%  Il nucleo del questionario — pagine fotocopiabili
%% =================================================================
%  Pagine da fotocopiare: niente spazi stirati a fondo pagina
%  (\raggedbottom vale da qui in poi; l'Appendice B lo imposta comunque)
\clearpage\raggedbottom
\section{Il nucleo del questionario}
\label{sec:indagine-nucleo}

{\small
\noindent\emph{Versione per l'università. Il questionario è anonimo:
nessuno potrà sapere che cosa hai risposto. Non ci sono risposte giuste
o sbagliate: interessa sapere che cosa fai davvero quando studi, non
che cosa \enquote{si dovrebbe} fare. Rispondi alle domande nell'ordine,
senza tornare indietro.}

\domanda{1.} Pensa a come ti prepari di solito a un esame. Elenca le
strategie che usi quando studi, \emph{in ordine di frequenza}: prima
quella che usi di più. Accanto a ciascuna scrivi perché la usi.

\begin{center}
  \begin{tabular}{@{}p{1.2em}|p{0.42\textwidth}|p{0.42\textwidth}@{}}
    & \textit{Strategia} & \textit{Perché la usi} \\
    \hline
    1 \rigavuota & & \\ \hline
    2 \rigavuota & & \\ \hline
    3 \rigavuota & & \\ \hline
    4 \rigavuota & & \\ \hline
    5 \rigavuota & & \\ \hline
    6 \rigavuota & & \\
  \end{tabular}
\end{center}

\domanda{2.} Di tutte le strategie che hai scritto, quale ti sembra la
più utile per imparare davvero? Perché?\par
\vspace{2.6\baselineskip}

\domanda{3.} Immagina di aver appena letto un capitolo del manuale per
un esame. Hai ancora un po' di tempo. Che cosa fai?
\begin{opzioni}
  \item Rileggo il capitolo, o una parte
  \item Provo a ricordare il contenuto senza guardare il libro (per
        esempio ripetendo a voce o scrivendo)
  \item Faccio qualcos'altro: \dotfill
\end{opzioni}

\domanda{4.} Se hai scelto di provare a ricordare: perché?
\begin{opzioni}
  \item Per capire quanto so e che cosa non so ancora
  \item Perché provare a ricordare mi aiuta a imparare e a ricordare
        meglio
  \item Perché è più veloce o più comodo
  \item Altro: \dotfill
\end{opzioni}

\domanda{5.} Qualcuno ti ha mai spiegato in modo esplicito come si
studia, cioè quali tecniche usare e perché?
\begin{opzioni}
  \item Sì, più volte \quad\quadretto\ Sì, una volta \quad\quadretto\ No
        \quad\quadretto\ Non ricordo
\end{opzioni}
Se sì, chi? (anche più risposte) \quadretto\ un insegnante o un
docente \quadretto\ un corso sul metodo di studio \quadretto\ i
genitori o altri familiari \quadretto\ un insegnante di ripetizioni
\quadretto\ video, social, siti \quadretto\ un'intelligenza
artificiale \quadretto\ un libro \quadretto\ compagni o amici

\domanda{6.} Pensa all'ultimo esame che hai sostenuto. Come hai
distribuito lo studio?
\begin{opzioni}
  \item Quasi tutto negli ultimi giorni (o nelle ultime notti)
  \item Soprattutto nell'ultima settimana o nelle ultime due
  \item Un po' ogni settimana durante il corso, con un ripasso finale
  \item In modo regolare fin dall'inizio, con riprese a distanza di
        tempo
\end{opzioni}

\domanda{7.} Quanto spesso usi ciascuna di queste strategie quando
studi? (una crocetta per riga)

\begin{center}\footnotesize\setlength{\tabcolsep}{2pt}\renewcommand{\arraystretch}{1.3}
  \begin{tabular}{@{}>{\raggedright\arraybackslash}p{0.46\textwidth}*{5}{>{\centering\arraybackslash}p{0.09\textwidth}}@{}}
    & \textit{mai} & \textit{raramente} & \textit{a volte} & \textit{spesso} & \textit{quasi\newline sempre} \\
    \midrule
    Rileggere il manuale o gli appunti & \quadretto & \quadretto & \quadretto & \quadretto & \quadretto \\
    Sottolineare o evidenziare & \quadretto & \quadretto & \quadretto & \quadretto & \quadretto \\
    Fare riassunti riscrivendo dal libro & \quadretto & \quadretto & \quadretto & \quadretto & \quadretto \\
    Fare riassunti o schemi a memoria, poi controllare & \quadretto & \quadretto & \quadretto & \quadretto & \quadretto \\
    Ripetere a voce senza guardare & \quadretto & \quadretto & \quadretto & \quadretto & \quadretto \\
    Farmi fare domande da qualcuno & \quadretto & \quadretto & \quadretto & \quadretto & \quadretto \\
    Usare flashcard (anche con un'applicazione) & \quadretto & \quadretto & \quadretto & \quadretto & \quadretto \\
    Rispondere a domande o rifare esercizi e vecchie prove & \quadretto & \quadretto & \quadretto & \quadretto & \quadretto \\
    Per questa riga scegli \enquote{spesso} & \quadretto & \quadretto & \quadretto & \quadretto & \quadretto \\
    Fare mappe guardando il libro & \quadretto & \quadretto & \quadretto & \quadretto & \quadretto \\
    Spiegare l'argomento a qualcun altro & \quadretto & \quadretto & \quadretto & \quadretto & \quadretto \\
    Chiedermi \enquote{perché?} e collegare a ciò che so & \quadretto & \quadretto & \quadretto & \quadretto & \quadretto \\
    Fare esercizi di tipi diversi mescolati & \quadretto & \quadretto & \quadretto & \quadretto & \quadretto \\
    Chiedere a un'IA di interrogarmi & \quadretto & \quadretto & \quadretto & \quadretto & \quadretto \\
    Chiedere a un'IA un riassunto o una risposta già fatta & \quadretto & \quadretto & \quadretto & \quadretto & \quadretto \\
    \bottomrule
  \end{tabular}
\end{center}

\domanda{8.} Leggi le due situazioni. Per ciascuna, scegli lo studente
che secondo te \emph{ricorderà di più dopo un mese}, e indica quanto sei
sicuro della tua risposta.

\emph{a)} Giulia e Marco hanno un'ora per studiare lo stesso capitolo.
Giulia lo legge quattro volte. Marco lo legge una volta; poi, per tre
volte, chiude il libro e scrive tutto ciò che ricorda, controllando dopo
ogni tentativo.\\
\quadretto\ Giulia \quad \quadretto\ Marco \quad \quadretto\ Più o meno
uguale \hfill Sicurezza (da 1 a 5): \ldots

\emph{b)} Sara studia tre ore il giorno prima dell'esame. Matteo studia
un'ora in tre giorni diversi della settimana prima.\\
\quadretto\ Sara \quad \quadretto\ Matteo \quad \quadretto\ Più o meno
uguale \hfill Sicurezza (da 1 a 5): \ldots

\domanda{9.} Su di te. Area del corso: \dotfill\ Anno di corso:
\ldots\ldots\ Almeno uno dei tuoi genitori è laureato? \quadretto\ sì
\quadretto\ no \quadretto\ non so

\medskip
\noindent\emph{Nella versione completa: altre situazioni di studio,
l'uso dell'intelligenza artificiale, il telefono e il sonno, le prove
intermedie, i motivi per cui studi, la media dei voti e, facoltativa,
la certificazione di DSA. Sulle risposte alle situazioni della domanda
8: capitoli~\ref{cap:recupero} e~\ref{cap:tempo}.}
\par}

\finecapitolo
